\subsection{Relation to local hidden-variable models}
  \label{subsec:lhv-relation}

  Violations of Bell inequalities are commonly interpreted as ruling out local
  hidden-variable models.
  More precisely, Bell’s theorem constrains a class of effective descriptions in which
  measurement outcomes are assumed to be determined by a set of underlying variables
  that renders joint probabilities factorizable for spacelike separated measurements.

  The analysis presented here clarifies that Bell inequalities do not directly constrain
  underlying physical dynamics.
  They constrain the logical structure of effective probabilistic descriptions.
  When observable outcomes admit an injective or effectively factorizable representation
  in terms of underlying variables, Bell-type inequalities apply.
  When this condition fails, Bell-type correlations may arise even if the underlying
  description is local, deterministic, and causally well behaved.

  Within the framework considered here, the failure of Bell inequalities originates
  from the non-injective nature of the mapping between underlying configurations and
  observable outcomes.
  Such non-injectivity arises generically when the effective description has a finite
  capacity to encode distinctions present at the underlying level.
  Because multiple underlying configurations correspond to the same effective
  description, no set of outcome-determining variables exists at the observable level
  that would allow joint probabilities to be decomposed into independent local
  contributions.

  Bell inequalities are therefore violated not because of nonlocal dynamics, but
  because the assumptions required for probabilistic factorization are not satisfied
  in finite-capacity effective descriptions.
